\begin{abstract}
Arbitrary-scale super-resolution (ASISR) predicts pixels at continuous coordinates, but recovery above the low-resolution Nyquist frequency and the role of scale-dependent readout remain difficult to assess. We study these questions with band-wise recovery measurements, covariance-conditioned linear references, and paired cell/footprint ablations. Across three datasets, normalising frequency by the sampling scale reduces variation in all nine evaluated model--dataset groups; the pre-specified collapse tolerance nevertheless passes only six. Evaluated decoders exhibit similar recovery in 10 of 12 low-reference-recovery cells, without establishing a universal decoder limit. Two further training seeds support coarse readout-loss ordering in all four architecture--seed groups. Averaging a full model over training-scale cell inputs predicts loss ordering, but fails function similarity in all four groups and improves on a simple footprint predictor in only two. Classical Wiener mismatch identities supply a conditional interpretation, not a network mechanism. Independent linear experiments corroborate approximate cost ranking, while a frozen 36-case guidance test makes a constant prior choice, trails cross-validation by .112 dB and costs 1.76 times as much. The resulting baseline is an empirical account of scale, frequency recovery and readout effects, with explicit failures of stronger mechanistic and adaptive-selection claims.
\end{abstract}
